% \iffalse meta-comment
%
% Copyright (C) 2017- by Yanshuo Chu <yanshuoc@gmail.com>
%
% This file may be distributed and/or modified under the
% conditions of the LaTeX Project Public License, either version 1.3a
% of this license or (at your option) any later version.
% The latest version of this license is in:
%
% http://www.latex-project.org/lppl.txt
%
% and version 1.3a or later is part of all distributions of LaTeX
% version 2004/10/01 or later.
%
% \fi
%
% \section{报告的子图兼容层}
% \changes{v3.2a}{2026/09/14}{报告的分图号补上括号、字号接进 \option{caption-sub}，照终稿}
%
%    \begin{macrocode}
%    \end{macrocode}
%
% ^^A ======================================================================
% ^^A Module report-subfig: 子图兼容层（hit-report 与 hitsz-interim-report 共用）
% ^^A ======================================================================
%    \begin{macrocode}
%<*report-subfig>
\bool_if:NF \g__hit_final_bool
  {
\hit@setup@subfigure
%    \end{macrocode}
%
% 子图题的字号与标签分隔是版式，留在类里；\cs{subcapsize} 共享层已经兜过底，
% 这里再探一次是因为报告可能自己定过。
%
% \emph{分图号要带括号。} 先前报告这一档没有重定义 \cs{thesubfigure}，
% \pkg{subcaption} 的默认值加上 \texttt{labelformat=simple} 排出来是光一个
% 「\texttt{a}」。本科指南 2.13.1 的正文写「（a）、（b）」，本科范例\emph{印出来}
% 是半角「\texttt{(a)~}」，终稿那一份早就按后者排了（见
% \file{src/flow/hit-final-float.dtx}），报告跟上。\cs{p@subfigure} 一并给，
% 交叉引用才排得出「图1-1 (a)」。
%
% \emph{字号也接进 \option{styles}。} 共享层的兜底是写死的 \cs{wuhao}；
% 终稿那一份读 \option{caption-sub} 目标，报告照做，用户写
% \option{styles}\texttt{=\{caption-sub=\{\ldots\}\}} 两档才是同一个口子。
%    \begin{macrocode}
\cs_set:Npn \thesubfigure { ( \alph { subfigure } ) }
\cs_set:Npn \thesubtable  { ( \alph { subtable } ) }
\cs_set:Npn \p@subfigure  { \thefigure \nobreakspace { } }
\cs_set:Npn \p@subtable   { \thetable \nobreakspace { } }
\cs_set:Npn \subcapsize
  { \docxlike_format_apply_font_else:nn { caption-sub } { \wuhao } }
\cs_if_exist:NF \subcapsize { \cs_new:Npn \subcapsize { \wuhao } }
\DeclareCaptionFont { hit@subcap } { \subcapsize }
\DeclareCaptionLabelSeparator { hit@sub } { \hit@subfigure@labelsep: }
\captionsetup [ sub ]
  { font = hit@subcap , labelformat = simple , labelsep = hit@sub , skip = 0pt }
  }
%</report-subfig>
%    \end{macrocode}
%
% \Finale
